% Options for packages loaded elsewhere
\PassOptionsToPackage{unicode}{hyperref}
\PassOptionsToPackage{hyphens}{url}
\PassOptionsToPackage{dvipsnames,svgnames,x11names}{xcolor}
\documentclass[
  11pt,
]{article}
\usepackage{xcolor}
\usepackage[margin=1in]{geometry}
\usepackage{amsmath,amssymb}
\setcounter{secnumdepth}{-\maxdimen} % remove section numbering
\usepackage{iftex}
\ifPDFTeX
  \usepackage[T1]{fontenc}
  \usepackage[utf8]{inputenc}
  \usepackage{textcomp} % provide euro and other symbols
\else % if luatex or xetex
  \usepackage{unicode-math} % this also loads fontspec
  \defaultfontfeatures{Scale=MatchLowercase}
  \defaultfontfeatures[\rmfamily]{Ligatures=TeX,Scale=1}
\fi
\usepackage{lmodern}
\ifPDFTeX\else
  % xetex/luatex font selection
  \setmainfont[Scale=1.0, BoldFont=texgyrepagella-bold.otf,
ItalicFont=texgyrepagella-italic.otf,
BoldItalicFont=texgyrepagella-bolditalic.otf]{texgyrepagella-regular.otf}
\fi
% Use upquote if available, for straight quotes in verbatim environments
\IfFileExists{upquote.sty}{\usepackage{upquote}}{}
\IfFileExists{microtype.sty}{% use microtype if available
  \usepackage[]{microtype}
  \UseMicrotypeSet[protrusion]{basicmath} % disable protrusion for tt fonts
}{}
\usepackage{setspace}
\makeatletter
\@ifundefined{KOMAClassName}{% if non-KOMA class
  \IfFileExists{parskip.sty}{%
    \usepackage{parskip}
  }{% else
    \setlength{\parindent}{0pt}
    \setlength{\parskip}{6pt plus 2pt minus 1pt}}
}{% if KOMA class
  \KOMAoptions{parskip=half}}
\makeatother
\usepackage{graphicx}
\makeatletter
\newsavebox\pandoc@box
\newcommand*\pandocbounded[1]{% scales image to fit in text height/width
  \sbox\pandoc@box{#1}%
  \Gscale@div\@tempa{\textheight}{\dimexpr\ht\pandoc@box+\dp\pandoc@box\relax}%
  \Gscale@div\@tempb{\linewidth}{\wd\pandoc@box}%
  \ifdim\@tempb\p@<\@tempa\p@\let\@tempa\@tempb\fi% select the smaller of both
  \ifdim\@tempa\p@<\p@\scalebox{\@tempa}{\usebox\pandoc@box}%
  \else\usebox{\pandoc@box}%
  \fi%
}
% Set default figure placement to htbp
\def\fps@figure{htbp}
\makeatother
% definitions for citeproc citations
\NewDocumentCommand\citeproctext{}{}
\NewDocumentCommand\citeproc{mm}{%
  \begingroup\def\citeproctext{#2}\cite{#1}\endgroup}
\makeatletter
 % allow citations to break across lines
 \let\@cite@ofmt\@firstofone
 % avoid brackets around text for \cite:
 \def\@biblabel#1{}
 \def\@cite#1#2{{#1\if@tempswa , #2\fi}}
\makeatother
\newlength{\cslhangindent}
\setlength{\cslhangindent}{1.5em}
\newlength{\csllabelwidth}
\setlength{\csllabelwidth}{3em}
\newenvironment{CSLReferences}[2] % #1 hanging-indent, #2 entry-spacing
 {\begin{list}{}{%
  \setlength{\itemindent}{0pt}
  \setlength{\leftmargin}{0pt}
  \setlength{\parsep}{0pt}
  % turn on hanging indent if param 1 is 1
  \ifodd #1
   \setlength{\leftmargin}{\cslhangindent}
   \setlength{\itemindent}{-1\cslhangindent}
  \fi
  % set entry spacing
  \setlength{\itemsep}{#2\baselineskip}}}
 {\end{list}}
\usepackage{calc}
\newcommand{\CSLBlock}[1]{\hfill\break\parbox[t]{\linewidth}{\strut\ignorespaces#1\strut}}
\newcommand{\CSLLeftMargin}[1]{\parbox[t]{\csllabelwidth}{\strut#1\strut}}
\newcommand{\CSLRightInline}[1]{\parbox[t]{\linewidth - \csllabelwidth}{\strut#1\strut}}
\newcommand{\CSLIndent}[1]{\hspace{\cslhangindent}#1}
\setlength{\emergencystretch}{3em} % prevent overfull lines
\providecommand{\tightlist}{%
  \setlength{\itemsep}{0pt}\setlength{\parskip}{0pt}}
\usepackage{xeCJK}
\setCJKmainfont{源雲明體月}
\usepackage{bookmark}
\IfFileExists{xurl.sty}{\usepackage{xurl}}{} % add URL line breaks if available
\urlstyle{same}
% fallback for those not using the hyperref driver hyperxmp:
\makeatletter
\@ifundefined{xmpquote}{\newcommand{\xmpquote}[1]{#1}}{}
\makeatother
\hypersetup{
  pdftitle={Research Statement},
  pdfauthor={Yu-You Liou},
  colorlinks=true,
  linkcolor={blue},
  filecolor={Maroon},
  citecolor={Blue},
  urlcolor={blue},
  pdfcreator={LaTeX via pandoc}}

\title{\textbf{Research Statement}}
\author{Yu-You Liou}
\date{2026-10-01}

\begin{document}
\maketitle

\setstretch{1.25}
{\raggedleft\small[\href{https://yyliou.github.io/yyliou/rs/rs_zh.pdf}{中文版由此去}]\par}

Air pollution, road deaths, and groundwater depletion share one feature.
Whoever causes the damage does not pay for it. Environmental and
resource economics calls this an externality, and policy is the
attempted correction. My research measures how much of that correction
arrives, and what else arrives with it. The answer is often not what the
policy assumed. I use causal inference and large-scale administrative
data. One line of work is on farms, the other on cities.

\subsection{1. Land, Food, and the
Environment}\label{land-food-and-the-environment}

Agriculture is the largest user of land and water, a major source of
emissions and nutrient runoff, and how millions of households live. A
policy that cleans up farming at the farmer's expense will not be
adopted. Must environmental quality and farm income trade off? Often
not.

Solar power is treated as a rival for farmland.
{[}\citeproc{ref-liou2025rooftop}{1}{]} and
{[}\citeproc{ref-lee2025solar}{2}{]} show that panels on existing farm
buildings displace nothing. Chicken farms that adopt them sell 5.8\%
more, from output rather than quality. The conflict is about siting, not
solar.

Siting can be set by policy. Taiwan zones land for aquavoltaics, which
places solar panels above working fish ponds instead of converting
farmland. {[}\citeproc{ref-chang2026aquavoltaics}{3}{]} finds that
farmland prices rose 13.5\% after land entered these zones. The increase
likely reflects the value of the development rights the zoning creates.

Fertilizer and pesticide runoff resists regulation because it comes from
tens of thousands of sources no inspector can reach.
{[}\citeproc{ref-chang2026food}{4}{]} finds that farms adopting food
traceability are 8.2, 15.4, and 40 percentage points less likely to use
chemical fertilizer, pesticide, and groundwater. Revenue rises too. A
food safety system did what enforcement cannot.

Biodiversity disappears from private land because nobody pays for it.
{[}\citeproc{ref-lee2025forest}{5}{]} finds that forest farms with
greater tree diversity earn less but face far steadier revenue, and
stability outweighs the loss. Diversity pays in stability, not income.
Programs that compensate forgone income misread the asset.

\subsection{2. Cities, Movement, and External
Costs}\label{cities-movement-and-external-costs}

Transport emissions keep rising, and transport kills directly. Cities
answer with rail, fare subsidies, and traffic rules. None is cheap, and
none has an obvious effect.

{[}\citeproc{ref-liou2026hookturn}{6}{]} uses the removal of the
two-stage left turn rule for motorcycles in part of Tainan City.
Accidents fell 21\%, victims 19\%, and vehicles involved 28\%. Road
safety usually means building something. Here a rule change cost nothing
and did as well.

Fare subsidies are now sold as climate policy.
{[}\citeproc{ref-liou2025monthly}{7}{]} finds that the Taipei monthly
pass raised ridership, passenger kilometers, and trip value by 4.3\%,
4.6\%, and 4.6\%. Emissions fall only if 70\% to 80\% of the new trips
come out of cars and motorcycles. Programs are judged on ridership.
Ridership is the wrong number. {[}\citeproc{ref-chang2021metro}{8}{]}
shows the demand is fragile as well, falling 1.43\% per COVID-19 case.

Mass events send the bill to people who never attended.
{[}\citeproc{ref-liou2026religious}{9}{]} finds accidents rose 14.9\%
and victims 17.4\% during the Dajia Mazu Pilgrimage, and 9.8\% even in
townships off the route. {[}\citeproc{ref-chang2024air}{10}{]} measures
the air quality cost of Hakka tomb sweeping at the Lantern Festival.
Neither cost enters the planning.

\subsection{3. Other}\label{other}

{[}\citeproc{ref-liou2025cbam}{11}{]} builds a strategic trade model
with green investment. The Carbon Border Adjustment Mechanism pushes
exporters toward cleaner technology, though its market structure effects
vary by industry. {[}\citeproc{ref-liou2024covid}{12}{]} finds that
consumers shifted food purchases online under pandemic risk.
{[}\citeproc{ref-liou2026hotel}{13}{]} finds that hotel revenue
recovered through price rather than volume, with restored labor supply
worth up to 17.8\%.

Rural or urban, one question runs through all of it. When policy changes
the price of environmental damage, who changes behavior, by how much,
and who pays.

\newpage

\subsection*{References}\label{references}
\addcontentsline{toc}{subsection}{References}

\protect\phantomsection\label{refs}
\begin{CSLReferences}{0}{0}
\bibitem[\citeproctext]{ref-liou2025rooftop}
\CSLLeftMargin{{[}1{]} }%
\CSLRightInline{Y.-Y. Liou, H.-H. Chang, and J. C. Begun. (2025).
{Rooftop Photovoltaics Adoption and Farm Revenue}. \emph{Agribusiness}.
\url{https://doi.org/10.1002/agr.70035}}

\bibitem[\citeproctext]{ref-lee2025solar}
\CSLLeftMargin{{[}2{]} }%
\CSLRightInline{T.-H. Lee, Y.-Y. Liou, and H.-H. Chang. (2025). {Is
Solar Panel Adoption a Win--Win Strategy for Chicken Farms? Evidence
from Agriculture Census Data}. \emph{Agriculture} 15, (20), 2124.
\url{https://doi.org/10.3390/agriculture15202124}}

\bibitem[\citeproctext]{ref-chang2026aquavoltaics}
\CSLLeftMargin{{[}3{]} }%
\CSLRightInline{H.-H. Chang, and Y.-Y. Liou. (2026). Does the
development of aquavoltaics policy affect farmland price? \emph{Land
Economics}. Conditionally accepted.}

\bibitem[\citeproctext]{ref-chang2026food}
\CSLLeftMargin{{[}4{]} }%
\CSLRightInline{H.-H. Chang, and Y.-Y. Liou. (2026). Food traceability
systems and the probability of using chemical fertilizers and pesticides
− empirical evidence using plot and farm household data from agriculture
census. \emph{Food Policy} 143, 103148.
\url{https://doi.org/10.1016/j.foodpol.2026.103148}}

\bibitem[\citeproctext]{ref-lee2025forest}
\CSLLeftMargin{{[}5{]} }%
\CSLRightInline{T.-H. Lee, Y.-Y. Liou, and H.-H. Chang. (2025). Forest
diversity and the distribution of farm revenue - {E}mpirical evidence
from forest farms in {T}aiwan. \emph{Forest Policy and Economics} 171,
103411. \url{https://doi.org/10.1016/j.forpol.2024.103411}}

\bibitem[\citeproctext]{ref-liou2026hookturn}
\CSLLeftMargin{{[}6{]} }%
\CSLRightInline{Y.-Y. Liou, and H.-H. Chang. (2026). The causal effects
of removing {Hook-Turn} regulation on road safety. \emph{Transportation
Research Part A: Policy and Practice} 205, 104860.
\url{https://doi.org/10.1016/j.tra.2026.104860}}

\bibitem[\citeproctext]{ref-liou2025monthly}
\CSLLeftMargin{{[}7{]} }%
\CSLRightInline{Y.-Y. Liou, M.-K. Kim, H.-H. Chang, and R.-W. Liou.
(2025). The causal effect of monthly pass programs on public
transportation and carbon emissions. \emph{Transportation Research Part
D: Transport and Environment} 146, 104903.
\url{https://doi.org/10.1016/j.trd.2025.104903}}

\bibitem[\citeproctext]{ref-chang2021metro}
\CSLLeftMargin{{[}8{]} }%
\CSLRightInline{H.-H. Chang, B. Lee, F.-A. Yang, and Y.-Y. Liou. (2021).
Does COVID-19 affect metro use in {T}aipei? \emph{Journal of Transport
Geography} 91, 102954.
\url{https://doi.org/10.1016/j.jtrangeo.2021.102954}}

\bibitem[\citeproctext]{ref-liou2026religious}
\CSLLeftMargin{{[}9{]} }%
\CSLRightInline{Y.-Y. Liou, and H.-H. Chang. (2026). Traffic accidents
of religious tourism. \emph{Annals of Tourism Research Empirical
Insights} 7, (1), 100209.
\url{https://doi.org/10.1016/j.annale.2026.100209}}

\bibitem[\citeproctext]{ref-chang2024air}
\CSLLeftMargin{{[}10{]} }%
\CSLRightInline{H.-H. Chang, Y.-Y. Liou, and C. D. Meyerhoefer. (2026).
{Air Pollution during Religious Holidays: Evidence from Tomb-Sweeping
Day}. \emph{Journal of the Association of Environmental and Resource
Economists}. Forthcoming. \url{https://doi.org/10.1086/743452}}

\bibitem[\citeproctext]{ref-liou2025cbam}
\CSLLeftMargin{{[}11{]} }%
\CSLRightInline{Y.-Y. Liou, H.-H. Chang, and J. C. Begun. (2025). The
impact of the carbon border adjustment mechanism policy on international
trade and market competition --- {A} theoretical justification.
\emph{Singapore Economic Review}.
\url{https://doi.org/10.1142/S0217590825500328}}

\bibitem[\citeproctext]{ref-liou2024covid}
\CSLLeftMargin{{[}12{]} }%
\CSLRightInline{Y.-Y. Liou, H.-H. Chang, and D. R. Just. (2024). How do
consumers respond to COVID-19? Application of {B}ayesian approach on
credit card transaction data. \emph{Quality \& Quantity} 58, (6),
5737--5754. \url{https://doi.org/10.1007/s11135-024-01915-9}}

\bibitem[\citeproctext]{ref-liou2026hotel}
\CSLLeftMargin{{[}13{]} }%
\CSLRightInline{Y.-Y. Liou, M.-K. Kim, and H.-H. Chang. (2026). Hotel
recovery strategies in the post-COVID-19 era: Evidence on revenue,
employment, and labor market rigidity from {T}aiwan. \emph{Journal of
Tourism Management Research} 13, (1), 148--162.
\url{https://doi.org/10.18488/31.v13i1.4940}}

\end{CSLReferences}

\end{document}
